\documentclass[bachelor, och, diploma]{SCWorks}
% Тип обучения (одно из значений):
%    bachelor   - бакалавриат (по умолчанию)
%    spec       - специальность
%    master     - магистратура
% Форма обучения (одно из значений):
%    och        - очное (по умолчанию)
%    zaoch      - заочное
% Тип работы (одно из значений):
%    coursework - курсовая работа (по умолчанию)
%    referat    - реферат
%  * otchet     - универсальный отчет
%  * nirjournal - журнал НИР
%  * digital    - итоговая работа для цифровой кафедры
%    diploma    - дипломная работа
%    pract      - отчет о научно-исследовательской работе
%    autoref    - автореферат выпускной работы
%    assignment - задание на выпускную квалификационную работу
%    review     - отзыв руководителя
%    critique   - рецензия на выпускную работу

% * Добавлены вручную. За вопросами к @mchernigin
\usepackage{preamble}


% Не оставлять одиночную строку абзаца на границе страниц.
\clubpenalty=10000
\widowpenalty=10000

\begin{document}

% Кафедра (в родительном падеже)
\chair{математической кибернетики и компьютерных наук}

% Тема работы
\title{Поиск изменений в связях между каналами мультивариативного временного ряда}

% Курс
\course{4}

% Группа
\group{451}

% Факультет (в родительном падеже) (по умолчанию "факультета КНиИТ")
% \department{факультета КНиИТ}

% Специальность/направление код - наименование
% \napravlenie{02.03.02 "--- Фундаментальная информатика и информационные технологии}
% \napravlenie{02.03.01 "--- Математическое обеспечение и администрирование информационных систем}
% \napravlenie{09.03.01 "--- Информатика и вычислительная техника}
\napravlenie{09.03.04 "--- Программная инженерия}
% \napravlenie{10.05.01 "--- Компьютерная безопасность}

% Для студентки. Для работы студента следующая команда не нужна.
% \studenttitle{студентки}

% Фамилия, имя, отчество в родительном падеже
\author{Карасева Вадима Дмитриевича}

% Заведующий кафедрой
\chtitle{доцент, к.\,ф.-м.\,н.}
\chname{С.\,В.\,Миронов}

% Руководитель ДПП ПП для цифровой кафедры (перекрывает заведующего кафедры)
% \chpretitle{
%     заведующий кафедрой математических основ информатики и олимпиадного\\
%     программирования на базе МАОУ <<Ф"=Т лицей №1>>
% }
% \chtitle{г. Саратов, к.\,ф.-м.\,н., доцент}
% \chname{Кондратова\, Ю.\,Н.}

% Научный руководитель (для реферата преподаватель проверяющий работу)
\satitle{доцент, к.\,ф.-м.\,н.} %должность, степень, звание
\saname{С.\,В.\,Папшев}

% Руководитель практики от организации (руководитель для цифровой кафедры)
\patitle{доцент, к.\,ф.-м.\,н.}
\paname{С.\,В.\,Миронов}

% Руководитель НИР
\nirtitle{доцент, к.\,п.\,н.} % степень, звание
\nirname{В.\,А.\,Векслер}

% Семестр (только для практики, для остальных типов работ не используется)
\term{4}

% Наименование практики (только для практики, для остальных типов работ не
% используется)
\practtype{учебная}

% Продолжительность практики (количество недель) (только для практики, для
% остальных типов работ не используется)
\duration{2}

% Даты начала и окончания практики (только для практики, для остальных типов
% работ не используется)
\practStart{01.07.2024}
\practFinish{13.01.2024}

% Год выполнения отчета
\date{2027}

\maketitle

% Включение нумерации рисунков, формул и таблиц по разделам (по умолчанию -
% нумерация сквозная) (допускается оба вида нумерации)
\secNumbering

\tableofcontents

% Раздел "Обозначения и сокращения". Может отсутствовать в работе
% \abbreviations
% \begin{description}
%     \item ... "--- ...
%     \item ... "--- ...
% \end{description}

% Раздел "Определения". Может отсутствовать в работе
% \definitions

% Раздел "Определения, обозначения и сокращения". Может отсутствовать в работе.
% Если присутствует, то заменяет собой разделы "Обозначения и сокращения" и
% "Определения"
% \defabbr

\intro

% После введения -- серии \section, \subsection и т.д.
Многомерные временные ряды возникают при одновременной регистрации нескольких показателей состояния одного объекта. К таким данным относятся измерения промышленных датчиков, телеметрия оборудования и показатели вычислительных систем. Каждый канал описывает отдельную характеристику, однако состояние объекта определяется не только значениями отдельных показателей, но и зависимостями между ними. Изменение режима работы может сопровождаться перестройкой этих зависимостей, даже если распределения значений отдельных каналов существенно не меняются.

Обнаружение моментов изменения временного ряда представляет собой задачу выявления переходов между участками с различными статистическими свойствами. В зависимости от постановки анализируется либо уже собранная запись, либо последовательность наблюдений по мере их поступления~\cite{truong2020review}. В настоящей работе рассматривается последовательное обнаружение изменений межканальных связей. При таком подходе решение в каждый момент времени должно основываться только на уже доступных измерениях.

Для представления зависимостей между каналами может использоваться взвешенный граф: вершины соответствуют каналам, а веса рёбер характеризуют силу выбранной статистической связи. Оценивание графа на последовательных временных окнах позволяет перейти от анализа отдельных сигналов к анализу изменения их совместного поведения. Задача поиска моментов перестройки корреляционных сетей рассмотрена, в частности, в работе И.~Барнетта и \mbox{Ю.-П.~Оннелы}~\cite{barnett2016correlation}. При этом статистическая связь между каналами сама по себе не устанавливает причинного влияния, а её изменение может отражать как нежелательное событие, так и допустимое переключение режима.

Одним из способов построения графа является вычисление коэффициента корреляции Пирсона, характеризующего линейную зависимость. Для зависимостей более общего вида может применяться взаимная информация, оценивание которой по конечной выборке рассматривается, например, в работе А.~Краскова и соавторов~\cite{kraskov2004mutual}. Выбор меры связи определяет, какие изменения способна обнаружить система. Вместе с тем более сложный оцениватель не гарантирует лучшего результата: его полезность необходимо проверять с учётом объёма данных в окне, шума и вычислительных затрат.

Основная трудность заключается в различении реальной перестройки зависимостей и случайных колебаний их выборочных оценок. Короткие окна позволяют быстрее учитывать новые наблюдения, но могут давать нестабильные оценки; длинные окна повышают устойчивость ценой задержки реакции. Поэтому актуальным является исследование методов обнаружения структурных изменений с одновременным учётом качества детекции, задержки и частоты ложных тревог.

Цель работы состоит в разработке и экспериментальном исследовании метода обнаружения изменений межканальных зависимостей многомерного временного ряда на основе динамики взвешенных графов связей, а также в определении условий его применимости по сравнению с анализом отдельных каналов.

Для достижения поставленной цели необходимо решить следующие задачи:
\begin{enumerate}
    \item Рассмотреть методы оценки межканальных зависимостей и обнаружения изменений в многомерных временных рядах.
    \item Создать синтетические данные с известными моментами смены зависимостей и изменившимися парами каналов; обеспечить воспроизводимость генерации и выбрать реальные данные для проверки обнаружения событий.
    \item Построить взвешенные графы связей в скользящих окнах, используя корреляцию Пирсона и взаимную информацию.
    \item Разработать правило обнаружения существенного изменения связей и исследовать способы уменьшения числа ложных тревог с учётом неопределённости оценок.
    \item Сравнить предложенный подход с методами анализа отдельных каналов и базовыми детекторами изменений при одинаковых условиях доступа к данным.
    \item Оценить качество обнаружения событий, задержку, частоту ложных тревог и вычислительную стоимость; исследовать влияние длины окна, шума и характера изменения, а на синтетических данных также проверить локализацию изменившихся связей.
\end{enumerate}

Объектом исследования являются многомерные временные ряды с изменяющимися зависимостями между каналами. Предмет исследования составляют методы оценивания этих зависимостей, построения динамических графов и обнаружения их структурных изменений.

Рабочая гипотеза состоит в том, что анализ динамики межканальных связей позволит обнаруживать часть изменений, слабо выраженных в статистических свойствах отдельных каналов. Предполагается также проверить, помогает ли учёт неопределённости оценок уменьшать число ложных срабатываний. Обе гипотезы требуют экспериментальной проверки при сопоставимых условиях сравнения.

Методическую основу работы составляют корреляционный и информационный анализ, представление зависимостей в виде графов, методы обнаружения смен и статистическое моделирование. Настройка параметров и порогов проводится на данных, отделённых от итоговой тестовой выборки. Качество обнаружения оценивается по событиям, а локализация изменившихся связей~--- на синтетических данных с известной структурой. Метки аномалий реальных наборов данных не рассматриваются как эталонная разметка изменения рёбер.

Практическая значимость работы заключается в возможности использования исследуемого подхода для мониторинга многоканальных данных, когда существенные изменения проявляются прежде всего в совместном поведении показателей. Планируемым результатом является программный прототип с воспроизводимыми экспериментами и рекомендациями по выбору меры связи, размера окна и правила обнаружения в зависимости от свойств данных.

\section{Теоретическая часть}

\subsection{Многомерные временные ряды и их свойства}
% TODO: Определения, обозначения, свойства рядов и пример многоканальных данных.

\subsection{Изменения значений и структуры зависимостей}
% TODO: Типы структурных смен, отличие от аномалий и последовательная постановка.

\subsection{Способы оценки межканальных зависимостей}
% TODO: Ковариация, корреляция Пирсона, взаимная информация, примеры и ограничения.

\subsection{Графовое представление зависимостей и временные окна}
% TODO: Матрицы связей, взвешенные графы, скользящие окна и интерпретация рёбер.

\subsection{Методы обнаружения изменений}
% TODO: Базовые детекторы, сравнение матриц, пороги, опорное состояние и CUSUM.

\subsection{Неопределённость оценок и устойчивость обнаружения}
% TODO: Ошибка оценки, длина окна, доверительные интервалы и блочный bootstrap.

\subsection{Критерии качества и правила сравнения}
% TODO: Событийные метрики, задержка, ложные тревоги, локализация и стоимость.

\subsection{Обзор исследований и обоснование выбранного подхода}
% TODO: Сравнение исследований, ограничения методов и исследовательские вопросы.

\clearpage

% Python 
% NumPy
% Pandas
% scikit-learn
% pymorphy2

% Отобразить все источники. Даже те, на которые нет ссылок.
% \nocite{*}

\bibliographystyle{ugost2003}
\bibliography{thesis}

%\begin{thebibliography}{99}

%\end{thebibliography}

% Окончание основного документа и начало приложений Каждая последующая секция
% документа будет являться приложением

\end{document}
